\documentclass[10pt,a4paper]{amsart}
\usepackage[margin=19mm]{geometry}
\usepackage{amsmath,amssymb,mathtools}
\usepackage[scr=rsfs,cal=euler]{mathalfa}
\usepackage{xfrac}
\usepackage[unicode,hidelinks,bookmarks=false]{hyperref}
\newcommand{\ie}{i.e.\ }
\newcommand{\cf}{cf.\ }
\newcommand{\resp}{respectively\ }
\newcommand{\Subsection}{\subsection}
\providecommand{\nobreakdash}{\nobreak\hbox{-}}
\setlength{\emergencystretch}{2em}
\multlinegap0pt
\usepackage{amssymb}
\usepackage{algorithm}\usepackage{algpseudocode}
\newtheorem{theorem}{Theorem}
\newcommand\xbf{\mathbf{x}}
\newcommand\ybf{\mathbf{y}}
\title{\LARGE \bf Learning-Based Stackelberg Equilibrium Seeking with Application to Demand-Side Energy Management}
\author{Silvia Cianchi, Reza Rahimi Baghbadorani, Anibal Sanjab, Sergio Grammatico}
\date{Source: arXiv:2605.00588v1 (CC BY 4.0)}
\begin{document}
\maketitle

\noindent\emph{Source and licence.} \LARGE \bf Learning-Based Stackelberg Equilibrium Seeking with Application to Demand-Side Energy Management. Version arXiv:2605.00588v1 (CC BY 4.0), distributed by arXiv under the Creative Commons Attribution 4.0 International licence, http://creativecommons.org/licenses/by/4.0/, as recorded in the arXiv licence metadata for this exact version and shown on the abstract page https://arxiv.org/abs/2605.00588v1. Full licence notice: \texttt{licenses/arxiv-2605.00588-CC-BY-4.0.txt}.\par\smallskip

\noindent\textit{Editorial scope.} Counted unit: the article's theorem Surrogate (source0.tex lines 279--290). The article's complete original proof of it is delivered with it (source0.tex lines 291--293, one \texttt{proof} environment of 3 lines). No mathematics is written, repaired or re-derived: the statement and the proof are the article's own lines, in the article's own notation, and no retained line is substituted or re-flowed.  No further source-local context is needed: every cross-reference of every retained line resolves inside the two delivered spans.  Nothing was omitted from any retained span, so the package carries no omission record.  No retained line contains a citation, so the wrapper carries no bibliography and no bibliography entry.  No retained line contains a figure environment, an image-inclusion command or drawing code, so no image is delivered and the package declares no asset dependency.  Editorial formatting only, in addition: an AMS \texttt{amsart} preamble that restates the article's own declarations and macros used by the retained lines, namely the environments \texttt{theorem} on separate counters, together with the standard packages those lines need.  The retained environments print in this document as Theorem 1 in order of appearance, whereas the article prints the same items with the numbering of its own full document.  The complete native source of the article as distributed in the arXiv e-print for this exact version is bundled unchanged as \texttt{sources/199577/source0.tex}.
\begin{theorem}\label{Surrogate}
    For each data point $(\ybf_k,\xbf_k) \in \mathcal{D}$, there exist $\epsilon>0, \, \rho >0$ such that:
    \begin{itemize}
        \item For all $\ybf \in \mathcal{B}(\ybf_k, \rho)$, it holds that \begin{align}
        \left\|\xbf^*(\ybf)-\xbf^*(\ybf; \hat{\Theta}_A, \hat{\Theta}_b)\right\|\le \epsilon;
        \end{align}
        \item Let $L$ be the Lipschitz constant of $J_0(\ybf,\xbf)$ with respect to $\xbf$, and $\hat{\ybf}^* \in \arg \min\limits_{\ybf \in \mathcal{B}(\ybf_k, \rho)}\, J_0(\ybf, \xbf^*(\ybf; \hat{\Theta}_A, \hat{\Theta}_b))$. Then it holds that 
        \begin{align}\label{jump}
        J_0(\hat{\ybf}^*, \xbf^*(\hat{\ybf}^*))\le J_0(\ybf_k, \xbf^*(\ybf_k))+ 2 L \epsilon.
        \end{align}
    \end{itemize}
\end{theorem}

\begin{proof}
    The proof can be derived by exploiting the fact that, at each data point, the suboptimality loss is equal to zero, and by combining the continuity of $x^*(\cdot)$, the Lipschitzness of $J_0$, and the definition of $\hat{y}^*$.
\end{proof}

\end{document}
